\documentclass[12pt, a4paper]{scrartcl}
\usepackage[lithuanian]{babel}
\setlength{\marginparwidth}{2cm}
\usepackage{../../1.\space Vienanariai\space ir\space daugianariai/manostilius_lt}
\renewcommand{\theproblem}{\arabic{problem}}
\usepackage{geometry}
\usepackage{tasks}
\settasks{label-width=2em, item-indent=3em}

\geometry{top=2cm, bottom=2cm, left=2cm, right=2cm}

\title{\textbf{Namų darbas: amžinybės dvelksmas}}
\author{Andrius Berniukevičius}
\date{2026-10-03}

\begin{document}

\maketitle

\section*{Namų darbo atlikimo tvarka}

\begin{itemize}[leftmargin=*]
    \item Užduotis atlikite A4 formato lapuose ir atneškite penktadienį. Darbai bus surinkti.
    \item Sprendimus rašykite įskaitomai ir tvarkingai. Neįskaitomas, netvarkingas arba pribraukytas darbas nebus priimtas – jį reikės perrašyti ir tą pačią dieną atnešti.
    \item Atsakymus pažymėkite spalva arba apibrėžkite stačiakampiu.
    \item Jei užstrigsite, klauskite per „Messenger“ arba „Teams“ – padėsiu.
    \item Nepalikite darbo paskutinei dienai.
\end{itemize}

\section{Sąlygos}

\begin{problem}
Apskaičiuokite nenaudodami smegenų minkštiklio:
\[
\frac{
    \left(\sqrt{2\sqrt{5}}-\sqrt{\sqrt{3}}\right)
    (2\sqrt{5}-\sqrt{3})
    \left(\sqrt{2\sqrt{5}}+\sqrt{\sqrt{3}}\right)
    +\sqrt{240}
}{
    \left(\sqrt[3]{3}-\sqrt[3]{2}\right)
    \left(\sqrt[3]{9}+\sqrt[3]{6}+\sqrt[3]{4}\right)
    \cdot \dfrac{\operatorname{MBK}(35,40)}{\operatorname{DBD}(840,360)}
}
\cdot \sqrt{\sqrt{\sqrt{7^8}}}
\]
\end{problem}

\begin{problem}
Įrodykite, jog žemiau esančio reiškinio reikšmė yra skaičius, lygus žuvų skaičiui, kurį apaštalas Petras su mokiniais pagavo po Kristaus nukryžiavimo. Daugyba stulpeliu ir smegenų protezo naudojimas draudžiamas:
\[
\frac{153^2-153\cdot 47}{212}
\cdot
\frac{145^3-15^3}{145\cdot 160+225}
\mathbin{:}
\frac{\sqrt{169}}{0{,}2}
\]
\end{problem}

\begin{problem}
Duota, jog $a$ yra lygties $x^2-2x-1=0$ sprendinys. Įrodykite, jog žemiau esančio reiškinio reikšmė lygi Evangelijų skaičiui:
\[
\frac{16a^4-4a^2}{8a^3+1}
\mathbin{:}
\frac{12a^3-3a}{12a^3-6a^2+3a}
\]
\end{problem}

\begin{problem}
Trijų iš eilės einančių teigiamų skaičių kvadratų suma lygi $194$. Įrodykite, jog šių skaičių vidurkis, pasuktas $90^\circ$ kampu, mūsų mintis nukreiptų į begalybę.
\end{problem}

\begin{problem}
Kvadrato viena priešingų kraštinių pora pailginta 5 vienetais, o kita pora sutrumpinta 3 vienetais. Gauto stačiakampio plotas lygus $20$. Apskaičiuokite nenaudodami prietaiso, dėl kurio žmogus lieka įkalintas matricoje:
\[
\sqrt{4P^2-144P+6^4},
\]
kur $P$ yra pradinio kvadrato perimetras.
\end{problem}

\begin{problem}
    Sudarykite redukuotas kvadratines lygtis, kurių sprendiniai yra šie skaičiai:
        \begin{tasks}[label={(\arabic*)}, label-offset=0.9em](3)
        \task $5$ ir $3$ 
        \task $-4$ ir $7$ 
        \task $-8$ ir $-5$
        \task $5+ \sqrt{7} $ ir $5- \sqrt{7}$
        \task $-3+ \sqrt{5} $ ir $-3- \sqrt{5}$
        \task $3\pi$ ir $-2\pi$  
    \end{tasks}
\end{problem}

\begin{problem}
    Stačiajame trikampyje statinių skirtumas lygus $3$, o įžambinė lygi $15$. Apskaičiuokite stataus trikampio plotą.
\end{problem}

\begin{problem}
    Išspręskite lygtis:
        \begin{tasks}[label={(\arabic*)}, label-offset=0.9em](2)
        \task $4 ^{x}+4 ^{-x} = \frac{17}{4}  $
        \task $x ^{4} = 5 x ^{2} -4 $
        \task $(x ^{2} -4x-3) ^{2} = (x+1)(x-5)+4$
        \task $\frac{x+1}{6-3 x}=\frac{x^2-4 x+4}{x^2+2 x+1} \cdot 2 \frac{2}{3}$    
    \end{tasks}
\end{problem}

\begin{problem}
    Išskaidykite dauginamaisiais kvadratinius trinarius:
    \begin{tasks}[label={(\arabic*)}, label-offset=0.9em](2)
        \task $4x^2 + 4x - 15$
        \task $5x^2 - 13x + 6$
        \task $6x^2 + x - 2$
        \task $4x^2 - 12x + 5$
        \task $3x^2 + 11x + 6$
        \task $5x^2 + 7x - 6$
        \task $6x^2 - 11x + 5$
        \task $8x^2 + 2x - 3$
    \end{tasks}
\end{problem}


\begin{problem}
    Žinoma, jog:
      \begin{tasks}[label={(\arabic*)}, label-offset=0.9em](1)
        \task lygties $x ^{2} +bx+8=0$ sprendinių skirtumas lygus $3,4$. Raskite $b$. 
        \task lygties $x ^{2} +8x+c=0$ sprendinių skirtumas lygus $2$. Raskite $c$.
        \task lygties $x ^{2} -12x+q=0$ sprendinių santykis $1:2$. Raskite $q$.
        \task lygties $x ^{2} +bx+15=0$ sprendinių santykis $3:20$. Raskite $b$.
    \end{tasks}
\end{problem}

\newpage
\section{Atsakymai}

\begin{enumerate}[label=\textbf{\arabic* užduotis.}, leftmargin=*]
    \item $69$
    \item ---
    \item ---
    \item ---
    \item $4$
    \item
    \begin{enumerate}[label={(\arabic*)}]
        \item $x^2-8x+15=0$
        \item $x^2-3x-28=0$
        \item $x^2+13x+40=0$
        \item $x^2-10x+18=0$
        \item $x^2+6x+4=0$
        \item $x^2-\pi x-6\pi^2=0$
    \end{enumerate}
    \item $54$
    \item
    \begin{enumerate}[label={(\arabic*)}]
        \item $x=-1$ arba $x=1$
        \item $x\in\{-2,-1,1,2\}$
        \item $x\in\{-1,\,2-\sqrt{6},\,2+\sqrt{6},\,5\}$
        \item $x=1$
    \end{enumerate}
    \item
    \begin{enumerate}[label={(\arabic*)}]
        \item $(2x-3)(2x+5)$
        \item $(x-2)(5x-3)$
        \item $(2x-1)(3x+2)$
        \item $(2x-5)(2x-1)$
        \item $(3x+2)(x+3)$
        \item $(5x-3)(x+2)$
        \item $(x-1)(6x-5)$
        \item $(2x-1)(4x+3)$
    \end{enumerate}
    \item
    \begin{enumerate}[label={(\arabic*)}]
        \item $b=\pm\frac{33}{5}$
        \item $c=15$
        \item $q=32$
        \item $b=\pm\frac{23}{2}$
    \end{enumerate}
\end{enumerate}

\end{document}
